\documentclass[12pt,a4paper]{article}

\usepackage{fontawesome5}
\providecommand{\faPencilAlt}{\faIcon{pencil-alt}}

\usepackage{styles/couleurs}
\usepackage{styles/mise_en_page}
\usepackage{styles/boites}
\usepackage{styles/maths}
\usepackage{styles/figures}

\usepackage{enumitem}
\usepackage{array}
\usepackage{tabularx}
\usepackage{amsmath}
\usepackage{amssymb}
\usepackage{tikz}
\usepackage{xparse}

\usetikzlibrary{calc,angles,quotes,arrows.meta}

\setlist[itemize]{leftmargin=1.4em,itemsep=0.15em,topsep=0.2em}
\setlist[enumerate]{leftmargin=*,itemsep=0.22em,topsep=0.25em}

\newcommand{\lignepointillee}{\noindent\dotfill\par\vspace{0.45em}}

\newcommand{\lignesreponse}[1]{%
\begin{tcolorbox}[
enhanced,colback=white,colframe=NeutreC!40,boxrule=0.5pt,arc=1.5mm,
left=3mm,right=3mm,top=2mm,bottom=1mm,before skip=2pt,after skip=2pt]
\foreach \i in {1,...,#1}{\lignepointillee}
\end{tcolorbox}}

\newcommand{\titreseq}[1]{%
\begin{center}
{\sffamily\LARGE\bfseries\color{TitreC}Séquence 6 -- Cosinus d'un angle aigu}\\[0.2em]
{\sffamily\large\bfseries\color{SavoirC}#1}
\end{center}}

\begin{document}

\titreseq{Activité 5 -- Calculer l'hypoténuse}

\section*{Partie A -- Comprendre l'opération inverse}

Dans le triangle \(JKV\) rectangle en \(K\) :

\[
JK=6{,}5\text{ cm}
\qquad\text{et}\qquad
\widehat{KJV}=28^\circ.
\]

On cherche \(JV\).

\begin{center}
\begin{tikzpicture}[scale=0.9]
  \coordinate (J) at (0,0);
  \coordinate (K) at (5.2,0);
  \coordinate (V) at (5.2,2.8);

  \draw[line width=0.9pt] (J)--(K)--(V)--cycle;
  \draw (4.82,0)--(4.82,0.38)--(5.2,0.38);

  \node[left] at (J) {$J$};
  \node[below] at (K) {$K$};
  \node[above] at (V) {$V$};

  \node[below] at (2.7,0) {$6{,}5$ cm};
  \node[above left] at (2.7,1.6) {$JV\,?$};

  \draw[ActionC,line width=1pt] (0.85,0) arc[start angle=0,end angle=28,radius=0.85];
  \node[ActionC] at (1.2,0.34) {$28^\circ$};
\end{tikzpicture}
\end{center}

Complète.

\[
\cos(28^\circ)=\frac{\ldots}{\ldots}
\]

Le côté adjacent est obtenu en multipliant l'hypoténuse par le cosinus :

\[
JK=JV\times\cos(28^\circ).
\]

Pour retrouver l'hypoténuse, on effectue l'opération inverse :

\[
JV=JK\div\cos(28^\circ).
\]

Donc :

\[
JV=6{,}5\div\cos(28^\circ)\approx\ldots\text{ cm}.
\]

Le résultat est-il bien supérieur à \(6{,}5\) cm ?

\section*{Partie B -- Méthode de rédaction}

\begin{tcolorbox}[colback=white,colframe=SavoirC,boxrule=0.7pt,arc=2mm]
Dans le triangle \(JKV\) rectangle en \(K\), on a :

\[
\cos(\widehat{KJV})=\frac{JK}{JV}.
\]

Donc :

\[
JV=JK\div\cos(\widehat{KJV}).
\]

Ainsi :

\[
JV=6{,}5\div\cos(28^\circ)
\approx\ldots\text{ cm}.
\]
\end{tcolorbox}

\section*{Partie C -- S'entraîner}

Calcule les hypoténuses demandées. Arrondis au dixième.

\[
\begin{array}{c|c|c|c}
\text{Triangle rectangle} & \text{Angle aigu} &
\text{Côté adjacent} & \text{Hypoténuse cherchée}\\
\hline
ABC\text{ rectangle en }B & \widehat{BAC}=41^\circ &
AB=7{,}2\text{ cm} & AC\\
\hline
DEF\text{ rectangle en }E & \widehat{EDF}=63^\circ &
DE=4{,}5\text{ cm} & DF\\
\hline
RST\text{ rectangle en }S & \widehat{TRS}=22^\circ &
RS=11\text{ cm} & RT
\end{array}
\]

\lignesreponse{10}

\section*{Partie D -- Analyser une erreur}

Lina cherche une hypoténuse. Elle écrit :

\[
6{,}5\times\cos(28^\circ)\approx5{,}7.
\]

Explique, sans refaire immédiatement tout le calcul, pourquoi son résultat
ne peut pas être l'hypoténuse.

Corrige ensuite son calcul.

\lignesreponse{5}

\aretenir{Pour calculer l'hypoténuse :

\[
\boxed{\text{hypoténuse}
=
\text{côté adjacent}\div\cos(\text{angle})}
\]

Diviser par un nombre compris entre \(0\) et \(1\) donne ici une longueur
plus grande que le côté adjacent.}

\end{document}
